% Noether P09 Korean producer draft — UNCHECKED
% T03-U19; ED0002 lines6554--6571; source 1,367 LF B / F3FE60232FBAAA7D16CCD7EE1D9D1943CEEB92E4E7B27D58D28097C5515BCEA6
% Pointer v009 / ED0002: B06BE3530D9CF2E82B56FDBA7FE41D5D044DF2425DFA2A059D4939EAA2F7A6C2 / C9A125167ACB33D914EE4374B65AE7CDF0052F568371B8B77B720EA178ABF0E3
% Translation only; no review, build, render, or approval.

$\mathfrak G$가 성질 I--IV를 만족한다는 증명은 체르멜로의 해당 증명과 정확히 나란히 진행되므로 간단히만 적겠다. 우선 $\mathfrak G$가 체르멜로 영역 $\mathfrak G_\eta$를 부분정역으로 포함하므로 II와 III은 명백하다. 또한 유리정수 계수 다항식에 관한 가우스 정리에 따르면 정수적 유리 함수식들의 합, 차, 곱도 다시 \emph{정수적} 유리 함수식이다. 따라서 $\mathfrak F$는 정역을 이룬다. 알려진 논법에 따라\srcfn{*)}{예컨대 Weber, § 97, 8 참조.} $\mathfrak G$도 정역이 되므로 I이 성립한다. 더 나아가 확장된 가우스 정리로부터\srcfn{**)}{Weber, § 20, 5.} 다음 두 보조정리가 따른다. ``$z$가 \emph{어떤} 방정식을 통해 $\mathfrak F$에 대하여 대수적으로 정수이면, 모든 유리 함수식의 체에 대하여 $z$가 만족하는 \emph{기약}방정식을 통해서도 그러하다.''\srcfn{***)}{Weber, § 97, 4 참조.} 그리고 ``어떤 대수적 수가 만족하는, 모든 유리수의 체에 대하여 기약인 방정식은 모든 유리 함수식의 체에 대해서도 기약이다.'' 그러므로 $\mathfrak G$는 분수형 대수수를 포함할 수 없다. 이로써 IV와 따라서 추상적 성질 I--IV가 모두 성립함이 증명되었다.
